%!TEX root=../../main.tex

\begin{lemma}\label{lemma:scaling-unique}
    Let \( \rbr{W^{(1)}, W^{(2)}} \) and \( \rbr{W^{(1)'}, W^{(2)'}} \)
    be two solutions to the non-convex ReLU training problem.
    If for every \( i \in [m]  \),
    it holds that
    \[
        W^{(1)}_{\cdot, i} W^{(2)}_{i} = W^{(1)'}_{\cdot, i} W^{(2)'}_{i},
    \]
    and \( \text{sign}(W^{(2)}_{i}) = \text{sign}(W^{(2)'}_{i}) \),
    then \( W^{(1)} = W^{(1)'} \) and \( W^{(2)} = W^{(2)'} \).
    That is, the solutions are the same.
\end{lemma}
\begin{proof}
    The ReLU prediction function \( f_{W^{(1)}, W^{(2)}} \) is invariant to
    scalings of the form
    \[
        \bar W^{(1)}_{\cdot, i} = \alpha W^{(1)}_{\cdot, i} \quad \bar W^{(2)}_{i} = W^{(2)}_{i} / \alpha,
    \]
    where \( \alpha > 0 \).
    Using this, we deduce that both solutions must satisfy the following
    equations:
    \begin{align*}
        1 & = \argmin_{\alpha} \alpha^2 \norm{W^{(1)}_{\cdot, i}}_2^2 + \norm{W^{(2)}_{i}}_2^2 / \alpha^2    \\
        1 & = \argmin_{\alpha} \alpha^2 \norm{W^{(1)'}_{\cdot, i}}_2^2 + \norm{W^{(2)'}_{i}}_2^2 / \alpha^2,
    \end{align*}
    which in turn implies that
    \[
        \norm{W^{(1)}_{\cdot, i}}_2 = \norm{W^{(2)}_{i}}_2, \quad \quad \norm{W^{(1)'}_{\cdot, i}}_2 = \norm{W^{(2)'}_{i}}_2
    \]
    We deduce
    \begin{align*}
        \norm{W^{(1)}_{\cdot, i}}^2_2
         & = \norm{W^{(1)}_{\cdot, i}}_2 \norm{W^{(2)}_{i}}_2    \\
         & = \norm{W^{(1)}_{\cdot, i} W^{(2)}_{i}}_2             \\
         & = \norm{W^{(1)'}_{\cdot, i} W^{(2)'}_{i}}_2           \\
         & = \norm{W^{(1)'}_{\cdot, i} }_2 \norm{W^{(2)'}_{i}}_2 \\
         & = \norm{W^{(1)'}_{\cdot, i}}^2_2,
    \end{align*}
    where we have used the fact that \( \norm{W^{(2)}_{i}}_2 = |W^{(2)}_{i}| \).
    This implies \( W^{(1)}_{\cdot, i} = W^{(1)'}_{\cdot, i} \) since \( W^{(1)}_{\cdot, i} \) and \( W^{(1)'}_{\cdot, i} \)
    are collinear, meaning \( W^{(2)}_{i} = W^{(2)'}_{i} \) must also hold.
\end{proof}

\begin{lemma}\label{lemma:mapping-equality}
    Let \( \rbr{W^{(1)}, W^{(2)}} \) and \( \rbr{\tilde W^{(1)}, \tilde W^{(2)}} \)
    be two solutions to the non-convex ReLU training problem.
    If \( \rbr{W^{(1)}, W^{(2)}} \) and \( \rbr{\tilde W^{(1)}, \tilde W^{(2)}} \) map
    to the same solution in the convex reformulation, then they are equal up to
    permutations or splits/merges of of neurons with collinear weights.
\end{lemma}
\begin{proof}
    Let \( (W^{(1)}, W^{(2)}) \) be a solution to the non-convex ReLU training
    problem.
    For each neuron, associate \( W^{(1)}_{\cdot, i} \) with the sparsest
    activation pattern \( D_i \) to which \( W^{(1)}_{\cdot, i} \) conforms.
    That is, given \( W^{(1)}_{\cdot, i} \), we associate it with the following pattern:
    \[
        D_i = \argmin \cbr{ \text{nnz}(D_j) : D_j \in \calD_Z, 2 (D_j - I) Z W^{(1)}_{\cdot, i} \geq 0 }.
    \]
    If \( W^{(1)}_{\cdot, i} \) and \( W^{(1)}_{\cdot, j} \) are associated with the same pattern
    \( D_i \) and \( \text{sign}(W^{(2)}_{i}) = \text{sign}(W^{(2)}_{j}) \),
    then \cref{chapter:fast-optimization} proves that \( W^{(1)}_{\cdot, i} \) and \( W^{(1)}_{\cdot, j} \)
    are collinear, i.e. \( W^{(1)}_{\cdot, i} = \alpha W^{(1)}_{\cdot, j} \) for some \( \alpha > 0 \).
    It also proves that the model obtained by merging these collinear neurons
    as
    \[
        \begin{aligned}
            W^{(1)'}_{\cdot, i}
             & = \frac{W^{(1)}_{\cdot, i}|W^{(2)}_{i}| + W^{(1)}_{\cdot, j}|W^{(2)}_{j}|}{\norm{W^{(1)}_{\cdot, i}|W^{(2)}_{i}| + W^{(1)}_{\cdot, j}|W^{(2)}_{j}|}_2^{1/2}}, \\
            W^{(2)'}_{i}
             & = \text{sign}(W^{(2)}_{i})
            \norm{W^{(1)}_{\cdot, i}|W^{(2)}_{i}| + W^{(1)}_{\cdot, j}|W^{(2)}_{j}|}_2^{1/2},
        \end{aligned}
    \]
    where \( W'_{1j} = 0 \), and \( W^{(2)}_{j} = 0 \),
    is also an optimal solution to the non-convex ReLU problem.
    We term this the merge/split symmetry.

    By recursively merging all collinear neurons assigned to the same
    activation pattern for which \( \text{sign}(W^{(2)}_{i}) =
    \text{sign}(W^{(2)}_{j}) \) holds, we eventually obtain an optimal
    \emph{reduced} model \( (W^{(1)'}, W^{(2)'}) \) such that
    no two non-zero neurons share an activation pattern \( D_i \) and have
    second-layer weights with the same sign.
    We may then place an ordering on the neurons by ordering the \( D_i \)
    matrices from \( 1 \) to \( p \).
    This removes all permutation symmetries.

    Since every signed neuron in the non-convex model now
    corresponds to a unique pattern \( D_i \), the mapping to the convex
    program is given by
    \[
        \begin{aligned}
            v_{i} & =
            \begin{cases}
                W^{(1)'}_{\cdot, i} W^{(2)'}_{i} & \mbox{if \( W^{(2)'}_{i} \geq 0 \)} \\
                0                                & \mbox{otherwise}
            \end{cases} \\
            u_{i} & =
            \begin{cases}
                -W^{(1)'}_{\cdot, i} W^{(2)'}_{i} & \mbox{if \( W^{(2)'}_{i} < 0 \)} \\
                0                                 & \mbox{otherwise},
            \end{cases}
        \end{aligned}
    \]
    Applying this procedure to \( \rbr{W^{(1)}, W^{(2)}} \) and \( \rbr{\tilde W^{(1)},
        \tilde W^{(2)}} \) and recalling that they map to the same solution implies
    \[
        W^{(1)'}_{\cdot, i} W^{(2)'}_{i} = \tilde W^{(1)'}_{\cdot, i} \tilde W^{(2)'}_{i},
    \]
    and \( \text{sign}(W^{(2)'}_{i}) = \text{sign}(\tilde W^{(2)'}_{i}) \)
    for every \( i \).
    Invoking \cref{lemma:scaling-unique} is sufficient to prove \( \rbr{W'_{1},
    w'_{2}} \) and \( \rbr{\tilde W^{(1)'}, \tilde W^{(2)'}} \) are identical.

    To obtain this result, we (i) merged all collinear neurons with the same
    sign in the second layer and then (ii) sorted the neurons according to
    their activation patterns.
    Thus, it must be that \( \rbr{W^{(1)}, W^{(2)}} \) and \( \rbr{\tilde W^{(1)}, \tilde W^{(2)}} \)
    differ only by (i) splitting of neurons into collinear neurons and (ii)
    the ordering (i.e. permutation) of the neurons.
    This completes the proof.
\end{proof}

\reluSolFn*
\begin{proof}
    Let the proposed optimal set be the completion of
    \[
        \begin{aligned}
            \calU_\lambda =
             & \big\{
            (W^{(1)}, W^{(2)}) :
            \, f_{W^{(1)}, W^{(2)}}(Z) =  \hat y,
            W^{(1)}_{\cdot, i} = (\sfrac{\alpha_{i}}{\lambda})^{\sfrac{1}{2}} v_i, \\
             & \hspace{1em} W^{(2)}_{i} = \xi_i (\alpha_i \lambda)^{\sfrac{1}{2}},
            \alpha_i \geq 0, \, i \in [2p] \setminus \calS_\lambda \Rightarrow \alpha_i = 0
            \big\},
        \end{aligned}
    \]
    over all permutations of the neurons and splits/merges of collinear
    neurons.
    Let \( (W^{(1)}, W^{(2)}) \in \calU_\lambda \).
    Inverting \cref{eq:convex-to-relu}, we compute a candidate solution \( u \)
    to the convex reformulation as
    \[
        u_i = W^{(1)}_{\cdot, i} W^{(2)}_{i} = \alpha_{i} v_i,
    \]
    for some \( \alpha_i \geq 0 \), where we assumed without loss of generality
    that \( W^{(2)}_{i} \geq 0 \) (if \( \xi_i = -1 \), then we simply map to a
    ``negative'' neuron in the convex reformulation).
    Moreover, \( \alpha_i = 0 \) if \( i \in [2p] \setminus \calS_\lambda \).
    Since the solution mapping preserves the prediction of the neural network,
    we also have \( X u = \hat y \), which is enough to guarantee \( u \)
    is a solution to the convex reformulation using \cref{prop:sol-fn-cgl}.
    Since \( u \) is a solution, we conclude that \( (W^{(1)}, W^{(2)}) \) must
    be optimal for the non-convex ReLU problem.

    For the reverse inclusion, let \( (W^{(1)}, W^{(2)}) \in \calO_\lambda \) and
    suppose no permutation or merge/split symmetry of \( (W^{(1)}, W^{(2)}) \)
    is in \( \calU_\lambda \).
    While \( (W^{(1)}, W^{(2)}) \not \in \calU_\lambda \), it still maps to
    some solution to the convex reformulation, call it \( u \).
    As \( \calU_\lambda \) is obtained by mapping every solution to the convex
    reformulation (i.e. every \( u \in \solfn_\lambda \) ) to the non-convex
    parameterization, it must be that \( u \) also maps to some
    \( (\tilde W^{(1)}, \tilde W^{(2)}) \) in \( \calU_\lambda \).
    \cref{lemma:mapping-equality} now shows that \( (W^{(1)}, W^{(2)}) \) and
    \( (\tilde W^{(1)}, \tilde W^{(2)}) \) are identical up to permutations and
    splits/merges of collinear neurons
    But, this implies \( (W^{(1)}, W^{(2)}) \in \calU_\lambda \) up to permutation
    and merge/split symmetries, which is a contradiction.
    We conclude that the optimal set is the completion of \( \calU_\lambda \) as
    claimed.
\end{proof}

\stationaryPoints*
\begin{proof}
    The proof is almost immediate.

    Given any sub-sampled set of activation patterns \( \tilde \calD \subset \calD_Z \),
    \citet[Theorem 3]{wang2022hidden} prove that the solutions to the sub-sampled
    convex program are Clarke stationary points \citep{clarke1990optimization} of
    the non-convex ReLU optimization problem in \cref{eq:non-convex-relu-mlp},
    and vice-versa.
    Using the expression for the CGL solution set in \cref{prop:sol-fn-cgl},
    which applies to sub-sampled convex reformulations as well as the full
    program, we obtain a version of \cref{cor:relu-sol-fn} for stationary
    points.
    That is, every model \( (W^{(1)}, W^{(2)}) \) in
    \begin{equation*}
        \begin{aligned}
            \calC_\lambda(\tilde \calD) = \big\{
             & (W^{(1)}, \! W^{(2)}) :
            \, f_{W^{(1)}, W^{(2)}}(Z) \! = \! \hat y_{\tilde \calD},
            W^{(1)}_{\cdot, i} = (\sfrac{\alpha_{i}}{\lambda})^{\sfrac{1}{2}} v_i(\tilde \calD), \,
            W^{(2)}_{i} = \xi_i (\alpha_i \lambda)^{\sfrac{1}{2}},                                      \\
             & \hspace{4em} \alpha_i \geq 0, \, i \in [m] \setminus \calS_\lambda \implies \alpha_i = 0
            \big\},
        \end{aligned}
    \end{equation*}
    is a stationary point of the non-convex ReLU program.
    Taking the union over all sub-sampled sets of activation patterns
    gives \( \calC_\lambda \), which is guaranteed to contain every stationary
    point of the non-convex objective.
    This completes the proof.
\end{proof}

\pUnique*
\begin{proof}
    Since there is only one solution to the convex reformulation, all solutions to the
    non-convex training problem must map to that solution.
    \cref{lemma:mapping-equality} now implies that the solution map for the
    non-convex problem is p-unique, i.e. unique up to permutations and
    splits/merges of collinear neurons.
\end{proof}

\reluUnique*
\begin{proof}
    We assume without loss of generality that only indices from \( 1 \)
    to \( p \) are in \( \equi \).
    By \cref{lemma:unique-cgl}, the constrained group lasso
    admits a unique solution
    if and only if
    \[
        \bigcup_{i \in \equi} \cbr{D_i Z \vi},
    \]
    are linearly independent.
    We now show that this fact holds under the proposed sufficient condition
    by proving the stronger fact that
    \( \bigcup_{i \in \equi} \cbr{[D_i Z]_j : j \in [d]} \)
    are linearly independent with probability one, where
    \( [D_i Z]_j \) is the \( j^{\text{th}} \)
    column of \( D_i Z \).
    Since \( \nnz(D_i) \geq d * p \) and \( Z \) has a continuous
    probability distribution, it holds that
    \( [D_i Z]_j \) has at least
    \( d * p \) non-zero entries with probability 1.
    Let
    \[
        \calS_{ij}
        = \text{Span}\rbr{\bigcup_{i \in \equi} \cbr{[D_i Z]_j : j \in [d]}
            \setminus [D_i Z]_j},
    \]
    and observe that \( \text{dim}(\calS_{ij}) \leq d * p - 1 \).
    As a result, the conditional probability \( [D_i Z]_j \) falls
    in this subspace satisfies
    \[
        \Pr([D_i Z]_j \in \calS_{ij} | \bigcup_{i \in \equi} \cbr{[D_i Z]_j : j \in [d]}
        \setminus [D_i Z]_j) = 0.
    \]
    Taking expectations over the remaining vectors in \( Z \) implies
    \[
        \Pr([D_i Z]_j \in \calS_{ij}) = 0.
    \]
    Finally, using a union bound over \( i, j \) implies that
    \( \bigcup_{i \in \equi} \cbr{[D_i Z]_j : j \in [d]} \) are linearly independent
    almost surely.
\end{proof}

\minimalComplexity*
\begin{proof}
    The proof follows directly from existing results.

    Recall from \citet{pilanci2020convexnn} that there are at most
    \[
        p \in O\rbr{r \frac{n}{r}^{3r}},
    \]
    activation patterns in the convex reformulation and that the complexity
    of computing an optimal ReLU model using a standard interior-point solver
    is \( O(d^3 r^3 (n / r)^{3r}) \).

    We know from \cref{prop:pruning-correctness} that the complexity of
    pruning an optimal neural network with at most \( 2p \) neuron
    is \( O(n^3 p + nd) \).
    Combining these complexities, we find that the cost of optimization dominates
    and overall complexity of computing an optimal and minimal neural network
    grows as \( O(d^3 r^3 (n / r)^{3r}) \).

    Finally, the bound on the number of active neurons follows from the
    fact that
    \[
        \dim \Span(\cbr{(X W^{*(1)}_{\cdot, i})_+}_i) \leq n.
    \]
    This completes the proof.
\end{proof}

\minNormNN*
\begin{proof}
    Let \( (u, v) \) be an optimal solution to the convex reformulation.
    \citet{pilanci2020convexnn} show that an optimal solution to the original
    two-layer ReLU optimization problem is given by setting
    \[
        W^{(1)}_{\cdot, i} = \frac{u_i}{\sqrt{\norm{u_i}_2}}, W^{(2)}_{i} = \sqrt{\norm{u_i}_2},
    \]
    and
    \[
        W^{(1)}_{\cdot, j} = \frac{v_i}{\sqrt{\norm{v_i}_2}}, W^{(2)}_{j} = -\sqrt{\norm{v_i}_2},
    \]
    where we define \( \frac{0}{0} = 0 \).
    Then, the \( r \)-value of any such solution can be calculated as
    \begin{align*}
        r(W^{(1)}, W^{(2)})
         & = \sum_{i = 1}^m \norm{W^{(1)}_{\cdot, i}}_2^4 + \norm{W^{(2)}_{i}}_2^4                                 \\
         & = \sum_{u_i \neq 0} \norm{\frac{u_i}{\sqrt{\norm{u_i}_2}}}_2^4 + \norm{\sqrt{\norm{u_i}_2}}_2^4
        + \sum_{v_i \neq 0} \norm{\frac{v_i}{\sqrt{\norm{v_i}_2}}}_2^4 + \norm{\sqrt{\norm{v_i}_2}}_2^4            \\
         & = \sum_{u_i \neq 0} \norm{u_i}_2^2 + \norm{u_i}_2^2 + \sum_{v_i \neq 0} \norm{v_i}_2^2 + \norm{v_i}^2_2 \\
         & = 2 \norm{u}_2^2 + 2 \norm{v}_2^2.
    \end{align*}
    That is, \( r(W^{(1)}, W^{(2)}) \) is a monotone transformation of the Euclidean
    norm of \( \rbr{u, v} \).
    Moreover, all optimal ReLU networks up to permutation and neuron-split
    symmetries can be obtained by applying this mapping to the solution to a convex reformulation.
    We conclude that the minimum \( r \)-valued optimal ReLU network before any
    such symmetries  are applied is given by the minimum \( \ell_2 \)-norm
    solution to the convex reformulation.

    Now, clearly permutation symmetries have no affect on the \( r \) value.
    However, splitting neurons into multiple collinear neurons can, in fact,
    reduce \( r \).
    For example, consider splitting a neuron \( (W^{(1)}_{\cdot, i}, W^{(2)}_{i}) \) into
    \( (W^{(1)'}_{\cdot, i}, W^{(2)'}_{i}) \) and \( (W^{(1)'}_{\cdot, j}, W^{(2)'}_{j}) \) such that
    \( W^{(1)'}_{\cdot, i} \) and \( W^{(1)'}_{\cdot, j} \) are collinear, \( \text{sign}(W^{(2)'}_{i}) =
    \text{sign}(W^{(2)'}_{j}) \), and
    \[
        \begin{aligned}
            W^{(1)}_{\cdot, i}
             & = \frac{W^{(1)'}_{\cdot, i} W^{(2)'}_{i} + W^{(1)'}_{\cdot, j} W^{(2)'}_{j}}{\norm{W^{(1)'}_{\cdot, i} W^{(2)'}_{i} + W^{(1)'}_{\cdot, j} W^{(2)'}_{j}}_2^{1/2}}, \\
            W^{(2)}_{i}
             & = \text{sign}(W^{(2)}_{i})\norm{W^{(1)'}_{\cdot, i} W^{(2)'}_{i} + W^{(1)'}_{\cdot, j} W^{(2)'}_{j}}_2^{1/2},
        \end{aligned}
    \]
    and \( \norm{W^{(1)'}_{\cdot, i}}_2 = \norm{W^{(2)'}_{i}}_2 \),
    \( \norm{W^{(1)'}_{\cdot, j}}_2 = \norm{W^{(2)'}_{j}}_2 \).
    As we showed in \cref{chapter:fast-optimization}, this
    operation preserves optimality \citep{mishkin2022convex}.
    However, the conditions on the split neurons also imply
    \begin{align*}
        \norm{W^{(1)}_{\cdot, i}}_2^4 + \norm{W^{(2)}_{i}}_2^4
         & = 2 \norm{W^{(1)'}_{\cdot, i} W^{(2)'}_{i} + W^{(1)'}_{\cdot, j} W^{(2)'}_{j}}_2^{2}                                   \\
         & = 2 \norm{W^{(1)'}_{\cdot, i} W^{(2)'}_{i}}_2^2 + 2 \norm{W^{(1)'}_{\cdot, j} W^{(2)'}_{j}}_2^2
        + 4 \abr{W^{(1)'}_{\cdot, i} W^{(2)'}_{i}, W^{(1)'}_{\cdot, j} W^{(2)'}_{j}}                                              \\
         & > 2 \norm{W^{(1)'}_{\cdot, i} W^{(2)'}_{i}}_2^2 + 2 \norm{W^{(1)'}_{\cdot, j} W^{(2)'}_{j}}_2^2                        \\
         & = \norm{W^{(1)'}_{\cdot, i}}_2^4 + \norm{W^{(2)'}_{i}}_2^4 + \norm{W^{(1)'}_{\cdot, j}}_2^4 + \norm{W^{(2)'}_{j}}_2^4,
    \end{align*}
    where have used positive collinearity of \( W^{(1)'}_{\cdot, i} \) and \( W_{2i}' \).
    As a result, the \( r \)-value of a solution can be decreased by splitting
    neurons.
    Thus, the minimum-norm solution to the convex reformulation is only the
    minimum \( r \)-value optimal ReLU network out of all
    networks which admit no neuron-merge symmetries.
\end{proof}

\oneNeuron*
\begin{proof}
    \citet{mishkin2022convex} show that \autoref{eq:non-convex-relu-mlp} has the same global optimal values as
    \begin{equation}\label{eq:mip}
        \min_{v, \gamma \in \cbr{-1, 1}} \half \norm{\sum_{i=1}^p (X v_i)_+ \gamma - y} + \lambda \sum_{i = 1}^p \norm{w_i}_2.
    \end{equation}
    Moreover, the objective-preserving mapping
    \( W^{(1)}_{\cdot, i} = v_i / \sqrt{\norm{v_i}_2} \),
    \( W^{(2)}_{i} = \gamma_i \sqrt{\norm{v_i}_2} \)
    can be used to obtain an optimal ReLU network from a solution to \cref{eq:mip}.
    We proceed by analyzing this equivalent problem and then use the
    mapping to return to the original non-convex formulation.

    Consider the case \( p = 1 \).
    Let \( X, y \) consist of two training points, \( (x_1, y_1) = (-100, 1) \) and \( (x_2, y_2) = (1, 10) \).
    In what follows, we drop the subscript for \( v \) and \( \gamma \) since \( p = 1 \) and we consider a one-dimensional.
    The optimization problem of interest is
    \begin{equation*}
        \min_{v, \gamma} \half \rbr{(x_1 v)_+ \gamma - y_1}^2 + \rbr{(x_2 v)_+ \gamma - y_2}^2 + \lambda \abs{v}.
    \end{equation*}
    Since \( x_1 < 0 \) and \( x_2 > 0 \), we can re-write this optimization problem as
    \[
        \min_{v, \gamma} \half \mathds{1}_{v \leq 0} \rbr{\rbr{x_1 v \gamma - y_1}^2 + y_2^2} + \mathds{1}_{v > 0}\rbr{\rbr{x_2 v \gamma - y_2}^2 + y_1^2} + \lambda \abs{v}.
    \]
    By inspection, we see that \( \gamma = +1 \) is optimal in both cases, leading to the following simplified expression:
    \[
        \min_{v} \half \mathds{1}_{v \leq 0} \rbr{\rbr{x_1 v - y_1}^2 + y_2^2} + \mathds{1}_{v > 0}\rbr{\rbr{x_2 v - y_2}^2 + y_1^2} + \lambda \abs{v}.
    \]
    This is a piece-wise continuous (but non-smooth) quadratic with a breakpoint at \( v^* = 0 \).
    We determine the solution to this minimization problem via a case analysis.

    \textbf{Case 1}: \( v^* = 0 \).
    Then, the optimal objective is trivially \( f(v^*) = y_1^2 + y_2^2 = 101 \).

    \textbf{Case 2}: \( v^* < 0 \).
    First order optimality conditions are
    \[
        x_1 \rbr{x_1 v_{-}^* - y_1} - \lambda = 0 \implies v_{-}^* = \frac{y_1 \cdot x_1 + \lambda}{x_1^2},
    \]
    which is valid only if \( \lambda < |y_1 \cdot x_1| = 100 \).
    The minimum objective value is then
    \begin{align*}
        \half \rbr{\rbr{x_1 v_{-}^* - y_1}^2 + y_2^2} - \lambda v_{-}^*
         & = \rbr{\rbr{x_1 \frac{y_1 \cdot x_1 + \lambda}{x_1^2} - y_1}^2 + y_2^2} - \lambda \rbr{\frac{y_1 \cdot x_1 + \lambda}{x_1^2}} \\
         & = \frac{\lambda^2}{x_1^2} + y_2^2 - \lambda \rbr{\frac{y_1 \cdot x_1 + \lambda}{x_1^2}}                                       \\
         & = -\frac{\lambda y_1}{x_1} + y_2^2                                                                                            \\
         & = \frac{\lambda}{100} + 100.
    \end{align*}

    \textbf{Case 3}: \( v_{+}^* > 0 \).
    Similarly to the previous case, we obtain
    \[
        x_2 \rbr{x_2 v_{+}^* - y_2} + \lambda = 0 \implies v_{+}^* = \frac{y_2 \cdot x_2 - \lambda}{x_2^2},
    \]
    which is valid only if \( \lambda < |y_2 \cdot x_2| = 10 \).
    In this case, the minimum objective is
    \begin{align*}
        \half \rbr{\rbr{x_2 v_{+}^* - y_2}^2 + y_1^2} + \lambda v_{+}^*
         & = \rbr{\rbr{x_2 \frac{y_2 \cdot x_2 + \lambda}{x_2^2} - y_2}^2 + y_1^2} - \lambda \rbr{\frac{y_2 \cdot x_2 + \lambda}{x_2^2}} \\
         & = \frac{\lambda^2}{x_2^2} + y_2^2 + \lambda \rbr{\frac{y_2 \cdot x_2 - \lambda}{x_2^2}}                                       \\
         & = \frac{\lambda y_2}{x_2} + y_1^2                                                                                             \\
         & = 10 \lambda + 1.
    \end{align*}

    To combine the cases, observe that
    \begin{align*}
        f(v_+^*) - f(v_-^*)
         & = \frac{\lambda y_2}{x_2} + y_1^2 - \sbr{-\frac{\lambda y_1}{x_1} + y_2^2} \\
         & = \lambda \rbr{\frac{x_1 y_2 + x_2 y_1}{x_1 x_2}} + y_1^2 - y_2^2          \\
         & = \lambda \rbr{\frac{-100 (10) + 1 (1)}{-100 (1)}} + 1 - 100               \\
         & = 9.99 \lambda - 99.
    \end{align*}
    We deduce that \( f(v_+^*) - f(v_-^*) > 0 \) (and thus \( v^* \leq 0 \))
    whenever \( \lambda > \frac{99}{9.99} \approx 10 \) and
    \( f(v_+^*) - f(v_-^*) > 0 < 0 \) otherwise.
    In this latter case, \( v^* \geq 0 \).

    Taking \( \lambda = 10 \), we find
    \[
        f(v_-^*) = 100 - 0.1 = 99.99 < 101 = y_1^2 + y_2^2 = f(0),
    \]
    so that \( v_-^* \) is optimal and \( v_-^* < 0 \).
    Moreover, \( v_-^* \) is strictly increasing as a function of \( \lambda \),
    for all \( \lambda > \frac{99}{9.99} \)
    so that \( v_-^* \) is optimal and strictly negative on the interval
    \( [\frac{99}{9.99}, 10] \).

    Now, consider \( \lambda = \frac{99}{9.99} - \epsilon \) to see that
    \[
        f(v_+^*) = \frac{990}{9.99} + 1 - 10 \epsilon < 101 = y_1^2 + y_2^2 = f(0),
    \]
    for all \( \epsilon > 0 \).
    We deduce that \( v_+^* \) is optimal for all \( \epsilon > 0 \) and
    thus \( v_+^* \) is optimal and strictly positive on the interval
    \( [0, \frac{99}{9.99}] \).

    To summarize, the solution function for this problem is as follows:
    \[
        \solfn(\lambda) =
        \begin{cases}
            \frac{\lambda}{100^2} - 0.01
             & \mbox{if \( \lambda > \frac{99}{9.99} \)}  \\
            \cbr{\frac{\lambda}{100^2} - 0.01, 0.1 - \frac{\lambda}{100}}
             & \mbox{if \( \lambda = \frac{99}{9.99} \)}  \\
            0.1 - \frac{\lambda}{100}
             & \mbox{if \( \lambda < \frac{99}{9.99} \).}
        \end{cases}
    \]
    This point-to-set map is clearly not open:
    for every sequence \( \lambda_k \uparrow \frac{99}{9.99} \),
    \( v_k \in \solfn(\lambda_k) \) implies
    \( \lim_k v_k \neq \frac{\lambda}{100^2} - 0.01 \).
    Moreover, a similar result holds for limits from above.
    Finally, it is clear by inspection that the optimal model fit is not
    unique at \( \lambda = \frac{99}{9.99} \), cannot be continuous
    in the functional sense, and, since it is not open, also fails to be
    continuous in the sense of point-to-set maps.
\end{proof}

